% ===========================================================================
% Supplementary Note S1 — exploratory Tier-1 land sector (LULUCF)
% Excluded from all agricultural totals/benchmarks/key-category/uncertainty in the main paper.
% ===========================================================================

\section*{S1.1 Overview and scope}

This Supplementary Note documents an exploratory Tier-1 land-use, land-use change and forestry (LULUCF)
extension of the agricultural inventory. It is \emph{excluded from every agricultural total, benchmark
comparison, key-category analysis, and uncertainty estimate} in the main paper; we report it here for
completeness because Pakistan's land sector is frequently cited in national climate discussion, but its
data basis is weaker than the agricultural inventory and its result should be read as indicative only.

\section*{S1.2 Method}

The land sector is computed separately from agriculture and is dominated by removals. Forest carbon is
estimated as a net afforestation sink consistent with the FAO Forest Resources Assessment forest-land
series \citep{fra2020}; Indus-delta mangrove blue carbon uses Global Mangrove Watch (GMW) area
\citep{gmw} and 2013 Wetlands Supplement carbon densities \citep{ipcc2013wetlands}; cropland and
grassland soil-carbon change are not estimated at Tier~1 (notation key NE, contributing nothing
to the net rather than a zero). The estimate is single-year.

\section*{S1.3 Result and its uncertainty}

A first-order Tier~1 calculation gives a small net sink of order \SI{-4}{}~Mt~\COe{} (forest plus
Indus-delta mangrove blue carbon). Both its magnitude and its \emph{sign} are uncertain: the
humid-tropical default carbon densities likely overstate storage in the arid delta, and although Pakistan
is often cited for mangrove afforestation, GMW\,v3 indicates a \emph{recent decline} in delta mangrove
extent, which would shrink or reverse the mangrove sink relative to the afforestation baseline. The land
factors used are summarised in Table~\ref{tab:s1land}; a properly developed land sector --- with
forest-area time series, stock-change factors, separated biomass and soil pools, and its own uncertainty
analysis --- is future work and is not attempted here.

\begin{table}[!htbp]
\centering
\footnotesize
\begin{tabular}{@{}llc@{}}
\toprule
\textbf{Land category} & \textbf{Basis} & \textbf{Notation} \\
\midrule
Forest land            & FAO FRA forest-land series & E \\
Cropland / grassland (SOC) & zero at Tier~1 & NE \\
Wetlands (mangrove)    & restored area $\times$ 6.5 t\COtwo{}/ha/yr (flux) & E \\
\bottomrule
\end{tabular}
\caption{Land-sector factors and notation keys (E~estimated, NE~not estimated). The
mangrove row is an annual FLUX---restored area at a sequestration rate---never area
$\times$ carbon density, which is a standing STOCK (the 471~\tCha{} humid-tropical
soil-carbon default puts the stock near 240~Mt~\COe{}, reported in the engine's
separate stock column and not as a removal). The stock--flux distinction is the reason
the note's sink is of order $-0.5$~Mt/yr for mangroves, not hundreds.}
\label{tab:s1land}
\end{table}

\section*{S1.4 Reproducibility}

The land sector is an optional module (\texttt{src/pkinv/lulucf/}) and is \emph{not} part of the default
agricultural reproduction run; it is skipped under \texttt{run\_all.py --quick} and its outputs
(\texttt{lulucf\_totals.csv}) feed only this note, never the agricultural totals or analyses of the main
paper.
